
\section{The Region-based Timing Tool} \label{sec:tool}
In this section, we dive into a powerful third-party performance tool built upon the \thiswork{} infrastructure that conducts region-based intra-kernel timing to assist in comprehending and optimizing the overlapping GPU hardware resources. 
To build a customized tool with \thiswork{}, the users must implement the instrumentation passes, the profile data management, and the post-processing process as introduced in \Sec{\ref{sec:third_party}}.
These are elaborated with the region-based timing tools as its workflow for instrumentation (\Sec{\ref{sec:profiler:workflow}}), its memory system (\Sec{\ref{sec:profiler:memory}}), and a trace replay post-processing (\Sec{\ref{sec:replay}}).


\subsection{Workflow} \label{sec:profiler:workflow}



We first introduce the workflow of the timing tool with two kinds of interfaces, the user interface and the compiler interface, as shown in \Fig{\ref{fig:workflow}}.
Depending on how the target program is parsed and the profiling regions are specified, different interfaces are invoked to instrument the profiling records.
\thiswork{} will then compile and execute the instrumented program to collect the raw performance counters.
At last, we replay the trace with the raw metadata and produce the final timeline trace.
This trace is output and feedback to different interfaces for visualization or profiler-guided compiler passes.


\paragraph{User interface} is the basic usage for developers to profile interested regions of a target GPU program.
We provide PythonDSL bindings in Triton kernels and also allow manual rewriting and overriding the dumped middle-level IR for the concerned regions.
It is straightforward as the developer manually rewrites the dumped middle-level IR for the concerned regions.
This is extremely useful when developing high-performance kernels and debugging efficiency problems.
The \thiswork{} will instrument the original program with rewritten program IR and produce the output trace.
Developers can utilize visualization tools, such as Chrome Tracer\cite{chrome_trace}, to get an intuitive view of the program timeline.

\paragraph{Compiler interface} is how the compiler interacts with the profiler.
Compared to the user interface, the compiler passes are responsible for parsing the middle-level IR and insert profiling record at concerned regions.
Similarly, the instrumented program is compiled and executed, and the profiled traces are feedback to the compiler passes within the compiler.

\subsection{Profile Data Management} \label{sec:profiler:memory}
We then introduce the specific management and decoding techniques of the tool.
The system accommodating the storage at each memory hierarchy is shown in \Fig{\ref{fig:design}}.
We demonstrate the data structure and layout bottom-up in the following.

\begin{figure}
    \centering
    \includegraphics[width=0.75\linewidth]{fig/design.pdf}
    \vspace{-0.4cm}
    \caption{Memory management of the region-based tool}
    \label{fig:design}
\end{figure}

\paragraph{Warp-level Region Data Structure.}
For each fine-grained profiling region, there are two profiling records inserted, the region start record and region end record.
Each record has 8 bytes, containing a 4-byte tag and a 4-byte payload, as shown in \Fig{\ref{fig:buffer}}.
The tag contains 1 flag bit as \emph{START} or \emph{END}, and 31-bit control bits that are dependent on the backend will be explained later. 
The payload contains a 32-bit time stamp metric captured from the hardware counter (cycle for \%clock in Nvidia and the LSB 32 bits of \code{S\_MEMTIME} instruction in AMD). 
This design achieves a good trade-off between space capacity and write speed. 
Because the storage consumption is low and the store takes a vectorized store instruction.

\R{
We use a 32-bit clock to capture the cycle tick of the current record, which may cause value overflow.
We address this in the post-processing procedure, where we detect and throw exceptions on overflows.
The trace replay can compensate for the clock wrap-around overflow as long as each iteration runs less than 4 billion cycles (4 seconds under 1GHz). 
Since most of the execution is spent on loops, the restriction is relaxed to the loop iteration level, which is less than 1 millisecond for most cases, making 32-bit cycle data a safe choice.
If necessary, \thiswork{} could easily extend to a 64-bit clock by adding dedicated operators and attributes in the IR.
}



\paragraph{SM-level Data Layout.}
The record is then copied to the data buffer in shared memory from the local register file with a store instruction.
The data buffer is split into profiling spaces by warp groups with non-overlapping record slots, as shown in \Fig{\ref{fig:design}}. 
As such, the record index at compile time determines the store location offset in shared memory.
For example, if we allocate 64 slots (0.5KB) and each thread block contains 2 warp groups (8 warps, 256 threads), then each warp group has 32 slots. 
During compile time, we pre-compute the base address of each profiling space and generate code to manage the indexes for each warp group.

The memory layout of each warp group buffer is shown in \Fig{\ref{fig:buffer}}. 
Start and end records are interleaved in shared memory to support nested regions and multiple iterations. 
The lower part of the figure illustrates three profiling patterns: common, nested, and multi-iteration. 
Records are stored without pairing and aligned during trace replay to ensure correctness based on the profiling patterns.

\paragraph{Shared Memory Circular Buffer.}
Because shared memory has limited capacity, it may not be able to accommodate all profile data, especially with a long multi-iteration record pattern.
For example, with 2 warp groups, 4 profiled regions, and a loop with 512 iterations, it takes up to  4096 record slots, translating to a 16KB shared memory storage. 
For the production-level kernels, we usually get the available shared memory capacity left ranging from 1KB to 4KB for intra-kernel profiling's data buffer, only up to 1.75\% of a total of 228KB on H100 device\cite{choquette2023nvidia}. 

\begin{figure}
    \centering
    \includegraphics[width=\linewidth]{fig/buffer.pdf}
    \vspace{-0.6cm}
    \caption{Circular buffer and record buffer}
    \label{fig:buffer}
\end{figure}



We propose a \emph{circular buffer} structure for each warp group's data buffer as a solution balancing accuracy and usability.
The circular buffer keeps only the trace's tail record cyclically when reaching the shared memory limit.
We overwrite the oldest results cyclically instead of flushing them to the lower-level memory.
The insight is that visualizing a few recent iterations is sufficient to identify the bottleneck.
We wrap around the index with a pre-calculated buffer capacity to implement the circular buffer.
If the captured events are overflowing the buffer, it is redirected to an occupied slot. 
This causes extra instructions for index management, which is addressed during compile-time.
We only need lightweight modular instructions to round the index at runtime.

\begin{figure}
    \centering
    \includegraphics[width=\linewidth]{fig/trace_replay.pdf}
    \vspace{-0.8cm}
    \caption{Trace replay}
    \vspace{-0.5cm}
    \label{fig:trace_replay}
\end{figure}

\R{
Besides, this \emph{circular buffer} is the default memory management behavior in the runtime.
The user can switch to the naive flush strategy, where the captured events are instantly written back to the global memory when the buffer is fully occupied.
This strategy can keep all the profile events with the overhead of many more frequent memory write instructions. 
In the \thiswork{}'s infrastructure, we provide a comprehensive operation to support the memory allocation at each hierarchy, including \code{Stack}, \code{Shared}, \code{Global}, and management strategies, including \code{Circular} and \code{Flush}.
Users can make use of these abstractions to develop their tool, combining the use of the buffers and management strategies.
}


\begin{figure*}
    \centering
    \includegraphics[width=\linewidth]{fig/trace_analysis.pdf}
    \vspace{-0.8cm}
    \caption{\OURR{Region-based timing results for FA3 kernels and overlapping improvements guided by profiling}}
    \vspace{-0.3cm}
    \label{fig:fa3_trace}
\end{figure*}

Specifically, we implement a collaborative store strategy for the AMD backend to minimize profiling overhead. 
Since only one timestamp record is profiled per warp or warp group, AMD GPUs utilize branching with thread masks to issue instructions with predication. 
However, this thread divergence can lead to unexpected instruction cache misses, causing overheads up to 600 cycles. 
To mitigate this, we enforce all threads to write to the same shared memory location, retaining only the last result. 
To align the records, we include an additional 12-bit signature derived from the least significant bits (LSB) of the \code{HARDWARE\_ID REG}, as shown in \Fig{\ref{fig:buffer}}. 
This signature encodes the \code{wave\_slot\_id}, \code{SIMD\_id}, and \code{pipe\_id} to annotate the profiled thread. 
At the warp group level, we ensure all four warps write to the same shared memory location.




\paragraph{Global Program Trace.}
Lastly, we gather the profiling data from all SMs back to the GPU global memory before returning from the kernel function. 
We keep the temporary profile data in the shared memory and copy the entire data buffer at once by the end of the kernel. 
Several extra metadata are attached for reconstructing the trace as shown in \Fig{\ref{fig:design}}, including thread block index, warp-level indices, and a number of recorded slots in each warp group. 
The Triton frontend automatically patches the kernel function with an extra argument (\code{profile\_mem}), allocating the GPU scratch memory for such profile data.


\subsection{Trace Replay} \label{sec:replay}



The major challenge in the region-based timing scenario of intra-kernel profiling is the perturbations caused by the profiling instructions.
Post-processing techniques, such as periodic sampling\cite{sample_based_profiling}, are usually used to mitigate the inaccuracies by sampling the performance metrics periodically and summarizing the results afterward to reduce the overhead associated with data collection.
\R{
In the region-based timing tool, we propose a similar post-processing method \emph{trace replay} to rescale the profiled regions according to the program semantics owing to the \thiswork{}'s design.
We can get accurate visualized results using this trace replay mechanism by offsetting a constant to account for the clock measuring overhead.
}





For regions with only synchronous instructions, we can easily rebuild the region by subtracting the overhead of the record operation.
For asynchronous instructions, we have two kinds of common inaccuracies as shown in \Fig{\ref{fig:trace_replay}}.
In the vanilla execution timeline example, we have a sequence of successive instructions executed after issuing an asynchronous MMA operation on the Tensor Core.
After the program reaches the barrier in the instruction stream, there is a period of waiting time before the functional unit can produce results.
To understand the asynchronous behavior here, users can insert profiling records in the successive execution to better align the operations.
However, the introduction of such records causes inaccurate wait time measurements.
In the \emph{reduced bubble} case, the wait time is underestimated due to the extra overhead caused by profiling.
Things are even worse in the \emph{unexpected idle} case, where the functional unit's execution time cannot cover the profiling overhead.
An unexpected idle time will be observed, and the optimization will be misleading.


\begin{table}[]
  \caption{Profiled regions of the FA3 kernel}
  \label{tbl:fa3}
  % \footnotesize
  \centering 
  \resizebox{\linewidth}{!}{
  \begin{tabular}{llll}
  \Xhline{3\arrayrulewidth}
  \multirow{2}{*}{Tag} & \multirow{2}{*}{Function}     & \multicolumn{2}{c}{Region ID} \\ \cline{3-4} 
                       &                               & Vanilla       & Improved      \\ \hline
  \rowcolor[HTML]{EFEFEF} 
  Load K               & Load the K tensor.            & 3      & 3       \\
  Load V               & Load the V tensor.            & 6      & 6       \\
  \rowcolor[HTML]{EFEFEF} 
  GEMM0                & Compute QK on TC.             & 12,22         & 12,22         \\
  GEMM1                & Compute PV on TC.             & 14,24         & 15,25         \\
  \rowcolor[HTML]{EFEFEF} 
  Softmax              & Compute softmax on CUDA core. & 13,23         & 17,27         \\ \Xhline{3\arrayrulewidth}
  \end{tabular}}
  \vspace{-0.5cm}
  \end{table}

To address this, we propose deducting the accurate wait time after profiling in the trace replay procedure as shown in \Fig{\ref{fig:trace_replay}}-(b).
Instead of placing one \code{END} record after the asynchronous instructions, we insert two \code{START} records before the asynchronous launch and after the wait barrier and one \code{END} record right before the barrier.
This guarantees an accurate wait time $T_{wait}$ measurement as $T_{wait} = (CLK_2-T_a) - (CLK_1-T_a)=CLK_2-CLK_1$, where $T_a$ and $T_b$ are the elapsed time in the record region before and after the clock read instruction respectively.
The profiling overhead is canceled out in the post-processing with the two carefully placed records.
We can derive the $T_{exe}$ time similarly.
This method requires the $T_{MMA}-T_{exe}>T_a+T_b$.
Since the profiling overhead is less than 25 cycles for most cases, as shown in \Sec{\ref{sec:evaluation:cycle}}, the execution time of the functional unit can cover only the record inserted here, which is usually about 1000 cycles.
